\documentclass[]{beamer}

\usepackage{graphicx}
\usepackage{booktabs}
\usetheme{Copenhagen}
\usefonttheme[onlymath]{serif}

\title{Parallelization of Diophantine Solvers}
\author{Hamza Iseric, Elijah Kin, Cameron Yeagle}
\date{May 5, 2026}

\begin{document}

\begin{frame}
  \titlepage
\end{frame}

\begin{frame}{Introduction}
  A \emph{Diophantine} equation is an equation for which we are only interested in integer solutions.
  \pause
  For instance, Pythagorean triples (e.g. $3^2 + 4^2 = 5^2$) are the solutions of the Diophantine equation $a^2 + b^2 = c^2$, while Fermat's last theorem famously states that $a^n + b^n = c^n$ has no solutions in the positive integers for $n \geq 3$.
  \pause
  \bigskip

  Question: Does
  \begin{equation*}
    a_1^6 + a_2^6 = b_1^6 + b_2^6 + b_3^6 + b_4^6 + b_5^6
  \end{equation*}
  have any solutions in the positive integers?
  \pause
  \bigskip

  Answer: Yes! The smallest was found in 2002:
  \begin{equation*}
    1117^6 + 770^6 = 1092^6 + 861^6 + 602^6 + 212^6 + 84^6
  \end{equation*}
\end{frame}

\begin{frame}{Resta-Meyrignac Algorithm}
  The crux of the search was that sixth powers mod $7$ can be congruent only to either $0$ or $1$ (and likewise for mod $8$ and mod $9$).
  \pause
  Then, if we are only interested in primitive solutions, at least one of $a_1$ and $a_2$ is not divisible by $7$.
  \pause
  Hence, mod $7$ the left-hand side is either $1$ or $2$.
  \pause
  Since there are only five terms on the right-hand side with each congruent to either $0$ or $1$, this means that at least three of the $b_i$ are divisible by $7$.
  \pause
  This yields the helpful restriction
  \begin{equation*}
    b_2^6 \equiv a_1^6 + a_2^6 - b_1^6 \mod 7^6.
  \end{equation*}
  \pause
  Finally, we define $v = a_1^6 + a_2^6 - (b_1^6 + b_2^6)$ and check if we can decompose $v / 7^6 = c_1^6 + c_2^6 + c_3^6$.
\end{frame}

\begin{frame}{Results}
  TODO Hamza
\end{frame}

\begin{frame}{Generalizability}

  New Equation:
  \begin{equation*}
    a_1^6 = b_1^6 + b_2^6 + b_3^6 + b_4^6 + b_5^6 +b^6_6
  \end{equation*}
  \pause
  Key Changes: 1 LHS term means $a_1^6$ is strictly congruent to $1$ modulo $7$. \pause This implies that only 1 RHS term is 1 modulo 7, rest are 0. Hence, 

  $$a_1^6\equiv b_1^6 \pmod{7^6}$$

  avoiding $O(N)$ search for $b_1$.
  
  
\end{frame}

\begin{frame}{Performance}

  \begin{figure}[h]
  \centering
  \resizebox{\textwidth}{!}{%
    \renewcommand{\arraystretch}{1.2}
    \begin{tabular}{lcccccccc}
      \toprule
      \textbf{\# Threads} & \textbf{1} & \textbf{2} & \textbf{4} & \textbf{8} & \textbf{16} & \textbf{32} & \textbf{64} & \textbf{128} \\
      \midrule
      \textbf{2 LHS, 5 RHS (ms)} & 136,298 & 71,504 & 36,329 & 18,402 & 9,186 & 4,842 & 2,618 & 1,466 \\
      \textbf{1 LHS, 6 RHS (ms)} & 127,541 & 67,311 & 47,525 & 29,703 & 20,703 & 11,160 & 9865 & 10745 \\
      \bottomrule
    \end{tabular}%
  } 
  \caption{Execution time scaling: 2 LHS, 5 RHS ($a_{max}=3000$) vs. 1 LHS, 6 RHS ($a_{max}=\text{6250}$).}
  \label{fig:runtime_scaling}
  \end{figure}

    
  \pause Difficulty: \textbf{Scalability}

  \begin{itemize}
      \item Inner loop must resolve 5 unknown variables
       
  \end{itemize}
  \bigskip
  
  No solutions found for $a_1<730,000$ (J.-C. Meyrignac, et al.). 
\end{frame}

\begin{frame}{References}
  \nocite{*}
  \bibliography{diophantine.bib}
  \bibliographystyle{abbrv}
\end{frame}

\end{document}